\chapter{Literature Review}
\label{chapter:literature-review}

% TODO - am I right in introducing the chapter in this fashion?
% is this the kind of pattern that will be repeated for all other chapters? i.e. introduce the chapter content, and talk about the upcoming sections and what to expect, etc.
In this chapter, there will be a review of the literature regarding the fundamentals of college course choice, found in Section \ref{section:fundamentals}.
As part of this review, a discussion of an individual's motivations in choosing a college course is provided in Section \ref{section:individual}.
Section \ref{section:external} seeks to explore the external and societal influences on an individual's college course choice.
Section \ref{section:theories-college} examines the theories of college and career choice in the realm of Vocational Psychology.
Similarly, Section \ref{section:theories-job} examines the theories of job and career satisfaction.
Interviews with expert guidance counsellors were conducted as part of this research, and a qualitative analysis of these interviews are provided in Section \ref{section:interviews}.
After a thorough review of the fundamentals of college course choice, Section \ref{section:relevant} seeks to answer the question: 
\emph{"What relevant information should we question a user in order to make suitable college course recommendations?"}.
In addition to this, this chapter provides an examination into closely-related work, found in Section \ref{section:closely-related-work}.
Finally, a summary of this chapter is provided in Section \ref{section:literature-review-summary}.

\section{Fundamentals of college course choice}
\label{section:fundamentals}

This section provides a broad overview into the fundamentals of college course choice.
Ultimately, it seeks to answer the question: \emph{"How do people pick college courses?"}.
As we will discover in the coming sections, there are no straightforward answers to this question, 
and there are many nuances which will be discussed.
Understanding the fundamentals of college course choice is important to this dissertation for two main reasons:
firstly, to help us answer the question: \emph{"What relevant information should the user be questioned in order to make suitable college course recommendations?"},
which will be explored in Section \ref{section:relevant}.
Secondly, to ensure that design decisions are well-informed by the literature surrounding college course choice, which will be explored in Chapter \ref{chapter:design}.

% TODO - mixed case or capitalise each word - what's the story here?
\subsection{Individual motivations}
\label{section:individual}

% i actually am not sure if it makes sense to split up these sections, could put them together,
% because sometimes an external influence is a reason for picking a course
% e.g. because parents wanted it
% e.g. because my friends/peers were doing it
% could also put college influences in this section here because it is relevant too

This section seeks to answer the question: \emph{"What reasons do people give for picking their college courses?"}.
As established in Chapter \ref{chapter:introduction}, choosing a college course is a major life decision,
and this section will examine the individual's motivations and reasons for picking their college courses.\\

% depends on gender, at least according to matthews2024students

\subsection{External influences}
\label{section:external}
% "what influences people in picking college courses?"

\subsection{College influences}
\label{section:influence}

% college itself
% prestige
% location
% financial limitations, can't move out - have to do a local college and stay at home
% sports
% friends

\subsection{Theories of career \& college course choice}
\label{section:theories-college}

% passion hypothesis

\subsection{Theories of job \& career satisfaction}
\label{section:theories-job}

\subsection{Interviews with guidance counsellors}
\label{section:interviews}

% TODO - is a summary necessary for a section? imo, yes
\subsection{Summary}
\label{section:fundamentals-summary}

\section{Relevant information to question}
\label{section:relevant}

% in light of the work done in section:fundamentals
% we must ask ourselves the question (which is a research objective): 
% what information is relevant to question a user in order to make effective, relevant recs for a CCRS?
% this is where the bias discussion comes in:
    % can cite my interesting paper where people prefer biased recs in the context of college courses
    % gender is established as a bias, how does that factor in?
    % parental education bg, SES, parental influence are also factors
    % i could make the argument that I wanted my RS to be "colour blind" to this

% don't want to overcomplicate and over-engineer the amount of preferences to ask a user

% TODO - do I talk about closely-related work first, then fundamentals, then relevant user information?
% TODO - should this be a separate chapter entirely?
\section{Closely-Related Work}
\label{section:closely-related-work}

% TODO - do I have a section for academic works?
% and then a separate section for enterprise products? e.g. FMCC

\section{Summary}
\label{section:literature-review-summary}

% Work in research areas tends to address a number of specific aspects. Ideally, the discussion of published research should focus on the aspects that have been addressed by various publications - and not a discussion of the individual publications.

% For example, if the topic would be a discussion of work on programming languages, the subsections of the related work could be discussions of object orientation and its realisation in various languages or the use of lambda functions by these languages.

